\DocumentMetadata{
  pdfversion=2.0,
  pdfstandard=ua-2,
  lang=en-US,
  tagging=on,
  tagging-setup={math/setup=mathml-SE,tabsorder=structure}
}
\documentclass[11pt,letterpaper]{article}
\usepackage[margin=0.68in,includefoot]{geometry}
\usepackage{newcomputermodern}
\usepackage{amsmath,amscd,mathtools}
\usepackage{tikz,adjustbox}
\providecommand{\square}{\mdlgwhtsquare}
\usepackage[hidelinks]{hyperref}
\hypersetup{
  pdftitle={Numerical Analysis PhD Exam, August 2019, Part 2},
  pdfauthor={University of Florida Department of Mathematics},
  pdfsubject={Graduate mathematics examination},
  pdfdisplaydoctitle=true,
  bookmarksopen=true
}
\setcounter{secnumdepth}{0}
\setlength{\parindent}{0pt}
\setlength{\parskip}{0.28em}
\setlength{\emergencystretch}{3em}
\setlength{\leftmargini}{2.15em}
\setlength{\leftmarginii}{2.65em}
\setlength{\leftmarginiii}{3em}
\renewcommand{\labelenumi}{\textbf{\theenumi.}}
\pagestyle{empty}
\raggedbottom
\allowdisplaybreaks
\newenvironment{examproblems}[1]{%
  \begin{enumerate}%
  \setcounter{enumi}{#1}%
  \setlength{\topsep}{0.35em}%
  \setlength{\partopsep}{0pt}%
  \setlength{\itemsep}{0.48em}%
  \setlength{\parsep}{0.18em}%
}{\end{enumerate}}
\newenvironment{examsubparts}{%
  \begin{enumerate}%
  \setlength{\topsep}{0.18em}%
  \setlength{\partopsep}{0pt}%
  \setlength{\itemsep}{0.16em}%
  \setlength{\parsep}{0.1em}%
}{\end{enumerate}}
\newenvironment{examinstructions}{%
  \par\begingroup\itshape
}{%
  \par\endgroup\vspace{0.18em}
}
\makeatletter
\renewcommand\section{\@startsection{section}{1}{0pt}%
  {-1.25ex plus -.25ex minus -.1ex}%
  {0.55ex plus .1ex}%
  {\normalfont\large\bfseries}}
\renewcommand{\@maketitle}{%
  \newpage\null\vskip -1em%
  {\footnotesize\bfseries
    \href{https://www.ufl.edu/}{University of Florida}, \href{https://math.ufl.edu/}{Department of Mathematics}\par}%
  \vskip -0.5em%
  \tagstructbegin{tag=Title}%
    {\LARGE\bfseries\href{https://gma.math.ufl.edu/exams/numerical-analysis-phd/}{Numerical Analysis PhD Exam}, August 2019, Part 2\par}%
  \tagstructend%
  \vskip 0.25em%
  \tagmcbegin{artifact}%
    \hrule height 1pt%
  \tagmcend%
  \vskip 1em%
}
\makeatother
\title{Numerical Analysis PhD Exam, August 2019, Part 2}
\author{}
\date{}
\begin{document}
\maketitle
\thispagestyle{empty}
\begin{examinstructions}
Do 4 (four) problems.
\end{examinstructions}
\begin{examproblems}{0}
\item
Consider using Newton’s method to find \(x \in \mathbb{R}^n\) such that \(F(x)=0\), for \(F:D\subseteq\mathbb{R}^n\to\mathbb{R}^n\), where~\(D\) is open and convex and~\(F\) is continuously differentiable on~\(D\).

\par
Suppose \(x^*\) satisfies \(F(x^*)=0\); and \(J(x)\), the Jacobian of~\(F\) evaluated at~\(x\), satisfies \(\|J^{-1}(x)\|\leq\mu\) for some number \(\mu>0\) for all~\(x\) in a convex neighborhood \(N\subseteq D\) that contains \(x^*\).
\begin{examsubparts}
\item[\textnormal{(a)}]
Assuming there is a constant \(0<\theta<1\) for which \(\|J(x)-J(y)\|\leq\theta/\mu\) for all \(x,y\in N\), show that \(\{x_k\}\) converges at least linearly to \(x^*\) whenever \(x_0\in N\).
\item[\textnormal{(b)}]
Assuming there is a constant~\(\kappa\) such that \(\|J(y)-J(x)\|\leq\kappa\|y-x\|\) for each \(x,y\in N\), show that \(\{x_k\}\) converges quadratically to \(x^*\) whenever \(x_0\in N\).
\end{examsubparts}
\item
Consider the data points \((x_1,y_1)=(0,-1)\), \((x_2,y_2)=(1,3)\), \((x_3,y_3)=(2,2)\).
\begin{examsubparts}
\item[\textnormal{(a)}]
Construct the both the Newton and Lagrange forms of the interpolating polynomial through \((x_1,y_1)\), \((x_2,y_2)\), \((x_3,y_3)\) (each can be left in the form of a linear combination of basis functions, but each basis function and coefficient should be explicitly shown).
\item[\textnormal{(b)}]
Let \(f(x)=1/x\) and show for \(x_0,\ldots,x_n\neq 0\) that 
\[
f[x_0,x_1,\ldots,x_n]=(-1)^n\prod_{i=0}^n\frac{1}{x_i}.
\]
\end{examsubparts}
\item
The third Chebyshev polynomial \(T_3\) is given by \(T_3(x)=4x^3-3x\).
\begin{examsubparts}
\item[\textnormal{(a)}]
Find the Chebyshev points on \(-2\leq x\leq 2\).
\item[\textnormal{(b)}]
Find the global interpolating polynomial on \(-2\leq x\leq 2\) that interpolates \(f(x)=e^x\) at the Chebyshev nodes.
\item[\textnormal{(c)}]
Suppose you were given 200 pieces of data, equally spaced on \(-2\leq x\leq 2\), and you want to approximate values that lie between data points. Explain what kind of interpolant you would use to fit these data, and why.
\end{examsubparts}
\item
Consider the function 
\[
s(x)=\begin{cases}
s_1(x)=(x+1)^3-3x, & -1\leq x\leq 0,\\
s_2(x)=(1-x)^3+3x, & 0\leq x\leq 1.
\end{cases}
\]
\begin{examsubparts}
\item[\textnormal{(a)}]
Is \(s(x)\) a cubic spline? If so, what kind?
\item[\textnormal{(b)}]
The trapezoid rule for numerical integration over \(a\leq x\leq b\) has an error given by 
\[
\int_a^b f(x)\,dx=I_T-\frac{1}{12}(b-a)^3f''(\eta),\qquad \text{for some }\eta\in(a,b).
\]
 Determine a bound for the error in the composite trapezoid rule assuming \([a,b]\) is broken up into~\(n\) equally spaced intervals of length \(h=(b-a)/n\).
\item[\textnormal{(c)}]
Use either the midpoint rule or the trapezoid rule with \(n=1\) then \(n=2\) to approximate \(\int_{-1}^1s(x)\,dx\). Explain which rule you chose to use and why.
\end{examsubparts}
\item
Let \(x_0=a\), \(x_1=a+h\) and \(x_2=b=a+2h\), and let \(f\in C^2[a,b]\).
\begin{examsubparts}
\item[\textnormal{(a)}]
Construct the difference approximation to \(f''(x_1)\) based on the \(P_2\) interpolant of~\(f\) on \([a,b]\) with interpolation points \(x_0,x_1,x_2\) (you should explicitly show how the difference approximation is derived from the interpolant).
\item[\textnormal{(b)}]
The \(p_2\) difference approximation satisfies \(\left|f''(x_1)-p_2''(x_1)\right|\leq Mh^2/12\), where~\(M\) is a constant that depends on~\(f\). Suppose the data is noisy and the approximation is based on the values \(f_i\) where \(f_i-f(x_i)=\epsilon_i\) with \(|\epsilon_i|<\epsilon\), \(i=0,1,2\), for a given value of~\(\epsilon\). What is the best accuracy with which \(f''(x_1)\) can be approximated? For what value of~\(h\) is it attained?
\end{examsubparts}
\end{examproblems}
\end{document}
